\subsection{\colorbox{blue!80}{\parbox{\textwidth}{PT12-PTH-TRV: Atak Path Traversal}}}

\subsubsection{Opis}
Atak path traversal polegał na próbie uzyskania nieautoryzowanego dostępu do systemu plików poprzez wykorzystanie mechanizmów przechodzenia po drzewie plików w systemie plików. Badania dotyczyły głównie końcówki /ftp/{ścieżka do pliku} i polegały na wstawianiu przejścia do katalogu wyżej "../" do podawanej ścieżki.

Atak w tej formie nie powiódł się, dlatego w szczegółach wykonania zawarto podjęte próby i wnioski z nich wynikające.

\subsubsection{Warunki wykonania ataku}
Atakujący musi odkryć ukrytą końcówkę ftp.

\subsubsection{Szczegóły wykonania}
Jeden z przykładów użycia "../" w został przedstawiony w tabelce. Przykład ten nie jest jeszcze skutecznym atakiem, ponieważ otrzymany plik wciąż należy do udostępnionego katalogu.
\begin{table}[H]
\begin{tabular}{|l|l|l|l|l|l|}
\hline
\rowcolor[HTML]{EFEFEF} 
\multicolumn{1}{|c|}{\cellcolor[HTML]{EFEFEF}\textbf{Metoda}} & \multicolumn{1}{c|}{\cellcolor[HTML]{EFEFEF}\textbf{Końcówka}} & \multicolumn{1}{c|}{\cellcolor[HTML]{EFEFEF}\textbf{Nagłowki}} & \multicolumn{1}{c|}{\cellcolor[HTML]{EFEFEF}\textbf{Ciasteczka}} & \multicolumn{1}{c|}{\cellcolor[HTML]{EFEFEF}\textbf{Dane}}                             & \multicolumn{1}{c|}{\cellcolor[HTML]{EFEFEF}\textbf{Odpowiedź/wynik}} \\ \hline
GET    & 
\begin{tabular}[c]{@{}l@{}}
/ftp/quarantine/../ \end{tabular} 
& -   & -   & - &  200 OK (przegląd plików dla /ftp) \\ \hline
\end{tabular}
\end{table}

W celu zwiększenia czytelności raportu reszta prób ataków została zawarta w raporcie w postaci listy złożonej ze ścieżki (końcówki) będącej celem żądania GET, odpowiedzi serwera oraz komentarza.

\textbf{WAŻNE!} Przy badaniach zaobserwowano, że wykonywanie zapytań przez przeglądarkę nie działa poprawnie, ponieważ stara się ona oszczędzić pracy serwerowi i skrócić ścieżkę jeszcze przed zapytaniem. Przykładem może być: \begin{lstlisting}
/ftp/quarantine/../ \end{lstlisting} Wysłanie takiego zapytania przez przeglądarke skutkowało zapytaniem bezpośrednim na /ftp widocznym w BurpSuite. Z tego powodu zapytania w tym przypadku musiały być wykonywane przez program BurpSuite.
\begin{itemize}
    \item \textbf{/ftp/../}\vspace{0.5em}
    
    ForbiddenError: Forbidden
    \item \textbf{/ftp/quarantine/../}\vspace{0.5em}
    
    listing directory /ftp/quarantine/../, czyli pliki /ftp/
    \item \textbf{/ftp/quarantine/\%2E\%2E/}\vspace{0.5em}
    
    listing directory /ftp/quarantine/.., czyli pliki /ftp/
    \item \textbf{/ftp/quarantine/\%252E\%252E}\vspace{0.5em}
    
    404 Error: ENOENT: no such file or directory, stat '/juice-shop/ftp/quarantine/\%2E\%2E'\vspace{0.5em}

    W tym przypadku można zauważyć, że do operacji dekodowania URL doszło tylko jeden raz (rozkodowany został znak \%, ale już nie kropka).
    \item \textbf{/ftp/\%2E\%2E}\vspace{0.5em}
    
    403 Error: Only .md and .pdf files are allowed!
    \item \textbf{/ftp/\%252E\%252E}\vspace{0.5em}

    403 Error: Only .md and .pdf files are allowed!\vspace{0.5em}

    Ten błąd jest tutaj zaskakujący. Prawdopodobnie jest on zwracany domyślnie tak, aby nie zdradzać zbyt wiele szczegółów użytkownikowi, jednak zalecane jest dokładniejsze spojrzenie na kod, czy aby na pewno program nie próbuje wykonać operacji w niedozwolonych lokalizacjach.
\end{itemize}

\subsubsection{Skutki}
Na tym etapie można podsumować działanie końcówki tak, że przechodzenie po folderach jest w pełni możliwe tak długo, jak użytkownik odwołuje się do czegoś wewnątrz ścieżki /ftp, czyli \textbf{atak path traversal nie jest wykonalny}. Zabezpieczenia są prawdopodobnie zrealizowane przez analizę przetworzonej ścieżki i sprawdzenie, czy jest wewnątrz /ftp, bądź przez odpowiednie ustawienia dostępu użytkownika (w kontekście użytkownika systemu kontenera serwera) reprezentującego serwer.

\subsection{\colorbox{yellow!80}{\parbox{\textwidth}{PT13-NL-BT: Atak Null Byte Poisoning}}}

\subsubsection{Opis}
Atak polegał na wstawieniu Null Byte w ścieżkę do pliku udostępnianego przez serwer w celu ominięcia ograniczenia możliwości pobierania plików do typów .md oraz .pdf. Dodanie null byte wraz z rozszerzeniem pdf pozwoliło na przejście ścieżki przez zabezpieczenie związane z typami, a serwer udostępniający pliki ignorował sztucznie dodane rozszerzenie i zwracał dowolny plik z serwisu ftp.

\subsubsection{Warunki wykonania ataku}
Atakujący musi znać końcówkę ftp oraz wiedzieć, że wykonanie ataku wymaga podwójnego zakodowania null byte'a. 

Dwustopniowe kodowanie URL, czyli użycie kodu \%2500 sprawiało, że w pierwszym etapie przetwarzania serwer rozkodowywał ścieżkę do \%00, a dopiero później, w drugim etapie, do null byte'a. Wiedza o tym pochodziła z zajęć laboratoryjnych. 

Atakujący musi również znać ścieżkę do innych plików w systemie (w przypadku testów wiedza ta pochodziła z ataku RCE).
\subsubsection{Szczegóły wykonania}
Przykład udanego ataku został zawarty w tabeli.
\begin{table}[H]
\begin{tabular}{|l|l|l|l|l|l|}
\hline
\rowcolor[HTML]{EFEFEF} 
\multicolumn{1}{|c|}{\cellcolor[HTML]{EFEFEF}\textbf{Metoda}} & \multicolumn{1}{c|}{\cellcolor[HTML]{EFEFEF}\textbf{Końcówka}} & \multicolumn{1}{c|}{\cellcolor[HTML]{EFEFEF}\textbf{Nagłowki}} & \multicolumn{1}{c|}{\cellcolor[HTML]{EFEFEF}\textbf{Ciasteczka}} & \multicolumn{1}{c|}{\cellcolor[HTML]{EFEFEF}\textbf{Dane}}                             & \multicolumn{1}{c|}{\cellcolor[HTML]{EFEFEF}\textbf{Odpowiedź/wynik}} \\ \hline
GET    & 
\begin{tabular}[c]{@{}l@{}}
/ftp/package.json.bak\%2500.md \end{tabular} 
& -   & -   & - &  200 OK (pobranie pliku package) \\ \hline
\end{tabular}
\end{table}

\subsubsection{Skutki}
Atak pozwala atakującemu na pobranie dowolnego pliku z serwisu ftp.

\subsubsection{Zalecenia}
Zabezpieczeniem przed udanym null byte poisoning była by przede wszystkim walidacja, czy plik wyszukany przez serwis udostępniający odpowiada temu, którego oczekiwała aplikacja sklepu. Oznacza to, że aplikacja powinna sprawdzić nie tylko czy ścieżka od użytkownika zawiera plik md/pdf, ale również czy to co chce zwrócić moduł udostępniający jest takim plikiem. Mogło by tutaj wystarczyć nawet podstawowe porównanie typu: pathFromUser == pathFoundByFtp.

\subsection{\colorbox{blue!80}{\parbox{\textwidth}{PT14-NL-BT-PTH-TRV: Atak path traversal z użyciem null byte poisoning}}}

\subsubsection{Opis}
Atak polegał na połączeniu opisanych wcześniej podejść path traversal oraz null byte poisoning w celu uzyskania dostępu do plików znajdujących się poza katalogiem udostępnianym przez ftp.

\subsubsection{Warunki wykonania ataku}
Atakujący musi znać ukrytą końcówkę ftp. Musi również znać ścieżkę do innych plików w systemie (w przypadku testów wiedza ta pochodziła z ataku RCE).

\subsubsection{Szczegóły wykonania}
Podjęte próby zostały zawarte w formie listy, gdzie każdy punkt zawiera końcówkę będącą celem zapytania GET oraz odpowiedź serwera.

W tym przypadku ważne jest użycie nadmiarowych przejść "wyżej" w drzewie (czyli '../'), aby upewnić się że ścieżka wyjdzie aż do katalogu-korzenia, czyli '/' 
\begin{itemize}
    \item \textbf{/ftp/../../../../../../../etc/hosts\%2500.md}\vspace{0.5em}
    
    400 Bad Request

    
    \item \textbf{/ftp/\%2E\%2E/\%2E\%2E/\%2E\%2E/\%2E\%2E/\%2E\%2E/\%2E\%2E/\%2E\%2E/\%2E\%2E/\newline
    \%2E\%2E/\%2E\%2E/\%2E\%2E/\%2E\%2E/\%2E\%2E/\%2E\%2E/\%2E\%2E/\%2E\%2E/etc/hosts\%2500.md}\vspace{0.5em}
    
    400 Bad Request
    \item \textbf{/ftp/\%252E\%252E/\%252E\%252E/\%252E\%252E/\%252E\%252E/\%252E\%252E/\newline
    \%252E\%252E/\%252E\%252E/\%252E\%252E/\%252E\%252E/\%252E\%252E/\%252E\%252E/\%252E\newline
    \%252E/\%252E\%252E/etc/hosts\%2500.md}\vspace{0.5em}
    
    403 Error: Only .md and .pdf files are allowed!
\end{itemize}

Opisana dokładniej została tylko końcówka ftp jednak warto wspomnieć, że istnieje drugie miejsce które może być podatne na path traversal:
\begin{table}[H]
\begin{tabular}{|l|l|l|l|l|l|}
\hline
\rowcolor[HTML]{EFEFEF} 
\multicolumn{1}{|c|}{\cellcolor[HTML]{EFEFEF}\textbf{Metoda}} & \multicolumn{1}{c|}{\cellcolor[HTML]{EFEFEF}\textbf{Końcówka}} & \multicolumn{1}{c|}{\cellcolor[HTML]{EFEFEF}\textbf{Nagłowki}} & \multicolumn{1}{c|}{\cellcolor[HTML]{EFEFEF}\textbf{Ciasteczka}} & \multicolumn{1}{c|}{\cellcolor[HTML]{EFEFEF}\textbf{Dane}}                             & \multicolumn{1}{c|}{\cellcolor[HTML]{EFEFEF}\textbf{Odpowiedź/wynik}} \\ \hline
GET    & 
\begin{tabular}[c]{@{}l@{}}
/assets/public/images/products/../../../../../../../../etc/hosts
\end{tabular} 
& -   & -   & - &  400 Bad Request \\ \hline
\end{tabular}
\end{table}

Innym przykładem jest zapytanie \lstinline[]|/assets/public/images/products/../products/apple_pressings.jpg|, które zwracało obrazek produktu, co pokazuje że chodzenie po ścieżce jest możliwe. W trakcie kilku prób testowych zachowanie nie różniło się końcówki /ftp/, dlatego ta końcówka nie została zbadana i opisana dokładniej.


\subsubsection{Skutki}
Nie udało się przeprowadzić skutecznego ataku, jednak zaobserwowano kilka faktów związanych z systemem.

W tym przypadku od razu w oczy rzuca się nowy błąd, czyli Bad Request. Biorąc pod uwagę odpowiedzi zwracane w badaniach Path Traversal (częste 403: only.md and .pdf [...] allowed) można sądzić, że przetestowane tutaj zapytania prawdopodobnie trafiają w plik /etc/hosts i z powodu Null Byte Poisoning serwer próbuje je pobrać, ale na pewnym etapie przetwarzania orientuje się, że jest to plik z niedozwolonej ścieżki i zwraca domyślny błąd. Jest to przesłanka, że obsługa błędów i zabezpieczenia mogą nie być zrealizowane dostatecznie profesjonalnie, dlatego zalecane było by przyjrzenie się działaniu tej końcówki od strony kodu.

\subsubsection{Zalecenia}
Mimo, że atak Path Traversal się nie powiódł to na podstawie zwracanych błędów można powiedzieć, że warto przyjrzeć się kodowi tej końcówki i upewnić się, że wszystko jest bezpieczne.

